% DRAFT: final numbers come from comparar_referencias.py on the final labor-only run.
\subsection{Comparison with published estimates}
\label{sec:benchmarks}

Table~\ref{tab:benchmarks} and Figure~\ref{fig:benchmarks} compare the estimate with every published annual series for Colombia that we could locate. Four families are available. The MIMIC estimates of Schneider and coauthors \citep{SchneiderBuehnMontenegro2010,MedinaSchneider2019} place the shadow economy between 35\% and 40\% of official GDP in the 1990s and early 2000s and between 25\% and 30\% after 2011. The enhanced MIMIC of \citet{AsllaniDellAnnoSchneider2024}, which restricts the concept to legal production deliberately concealed from the authorities, is almost flat at 24\%--27\%. The World Bank database \citep{ElginKoseOhnsorgeYu2021} reports a DGE series, built on the model of \citet{ElginOztunali2012}, that declines from 35\% in 1991 to 29\% in 2020, and a MIMIC series of 34\%--37\%. \citet{Vuletin2008} estimates 43.5\% for the early 2000s.

Published estimates differ among themselves by up to {{SPREAD}} percentage points in the same year. Our baseline lies inside that range throughout: {{BASE_VS_REFS}}. Its decline is steeper than that of the MIMIC series, which change little by construction, and similar to that of the DGE series of the World Bank, whose informal technology is the same as ours. The equal-technology specification is above every published series except \citet{Vuletin2008}, and the aggregates-only estimate is below all of them.

Three differences in concept qualify the comparison. First, MIMIC estimates are expressed as a percentage of official GDP and assume that the shadow economy is not in it, whereas Colombian national accounts already include an estimate of the non-observed economy; our share is a fraction of measured GDP. Second, MIMIC levels are not independent measurements: they are converted to percentages of GDP with benchmark values from other studies, and \citet{AsllaniDellAnnoSchneider2024} anchor the scale to informal employment multiplied by an assumed informal wage, which is the accounting of Section~\ref{sec:model} with relative productivity assumed instead of measured. Third, our concept is the output of informal workers, which includes legal own-account production that is not concealed and excludes illegal activity.

\begin{figure}[H]
\centering
\includegraphics[width=0.9\textwidth]{fig_benchmarks.pdf}
\caption{Informal share of output: this paper and published estimates for Colombia (percent of GDP).}
\label{fig:benchmarks}
\end{figure}
